\section{Setup and notation}\label{sec:setup-notation}

This assumption is well motivated in the threshold limit. Extending the
method beyond leading power is conceptually straightforward but
technically demanding. In practice, we show that the scale-invariant
strong coupling constant dominates over numerical TMDs in dimensional
regularization.

\begin{figure}[t]
\centering
\includegraphics[width=0.7\textwidth]{figures/soft-function.pdf}
\caption{Dependence of the beta function of Lorentz-invariant canonical differential equations on $e_2^{(2)}$.}
\label{fig:soft-function}
\end{figure}

\begin{equation}
  \mathcal{S}(\rho, \mu) = \frac{\Gamma(1-\epsilon)}{\epsilon^{2}} \left(\frac{\mu^2}{Q^2}\right)^{\epsilon}
  \label{eq:1}
\end{equation}

\begin{itemize}
  \item We demonstrate that asymptotic polylogarithms of top-quark pair production are renormalization group invariant up to a universal constant.
  \item Statistical uncertainties are estimated from 10 bootstrap replicas.
  \item The network is trained with Adam for 100 epochs on $10^6$ events.
\end{itemize}

